\chapter[Eigenvalues \& Diagonalization]{Polynomial Regression, Eigenvalues, Eigenspaces \& Diagonalization}
\label{chap:week04}

% ----------------------------------------------------------------------
% PART 1: EXECUTIVE ROADMAP & LEARNING OBJECTIVES
% ----------------------------------------------------------------------
\section{Executive Roadmap \& Learning Objectives}

At first glance, linear models appear fundamentally limited: they can only construct flat hyperplanes and linear decision boundaries. Yet, in modern data science, linear algorithms are routinely deployed to capture highly non-linear physical phenomena, aerodynamic drag curves, and financial volatility smiles. This remarkable versatility is unlocked through \textbf{non-linear feature expansions}: by projecting low-dimensional inputs into higher-dimensional feature spaces via polynomial basis functions, a linear model learns complex curves while retaining the convex optimization guarantees of Ordinary Least Squares.

However, once we operate in expanded feature spaces or analyze iterated transformations (such as recurrence relations in deep networks or transitions in Markov models), we confront matrices acting repeatedly on vector states: $\mathbf{x}_{k} = A^k \mathbf{x}_0$. Computing $A^k$ directly via brute-force matrix multiplication costs $\mathcal{O}(k n^3)$ operations and tangles coordinate axes uncontrollably. 

The mathematical remedy is the \textbf{spectral theory of linear operators}. By identifying the invariant directions—the \textbf{eigenvectors}—and their corresponding scale factors—the \textbf{eigenvalues}—we decouple a coupled matrix transformation into $n$ independent scalar scalings. This chapter builds the theoretical and practical bridge from polynomial feature representations to the spectral decomposition of linear operators.

\subsection{Prerequisites \& Scaffolding}
Mastery of this chapter builds upon vector spaces, determinant calculations, and ordinary least squares, as outlined in Table~\ref{tab:w04_prereqs}.

\begin{table}[htbp]
\centering
\small
\begin{tabularx}{\textwidth}{lXX}
\toprule
\textbf{Foundational Area} & \textbf{Core Concepts Required} & \textbf{Role in Machine Learning Foundations} \\
\midrule
Linear Algebra & Determinants $\det(A)$, matrix invertibility, linear independence & Solving characteristic polynomial $\det(A - \lambda I) = 0$ \\
Subspace Theory & Nullspace $\Null(A - \lambda I)$, rank-nullity theorem & Computing geometric eigenspaces and diagnosing deficiency \\
Least Squares & Design matrix $X$, normal equations $(X^T X)\mathbf{w} = X^T \mathbf{y}$ & Fitting polynomial curves via Vandermonde expansions \\
\bottomrule
\end{tabularx}
\caption{Prerequisite Scaffolding for Chapter 4.}
\label{tab:w04_prereqs}
\end{table}

\subsection{Core Competencies}
By the conclusion of this chapter, students will be able to:
\begin{itemize}
    \item Formulate polynomial regression problems via the non-linear feature map $\boldsymbol{\phi}(x)$ and construct Vandermonde design matrices.
    \item Compute eigenvalues, eigenvectors, and geometric eigenspaces $E_\lambda = \Null(A - \lambda I)$ for arbitrary square matrices.
    \item Distinguish rigorously between algebraic multiplicity $am(\lambda)$ and geometric multiplicity $gm(\lambda)$ to diagnose defective matrices.
    \item Diagonalize square matrices into $A = PDP^{-1}$ and compute iterated matrix powers $A^k = P D^k P^{-1}$ in closed form.
    \item Formulate discrete linear dynamical systems (including recurrence relations and Markov chains) and determine their asymptotic stability.
\end{itemize}

\subsection{Notation \& Mathematical Conventions}
\begin{itemize}
    \item $A \in \mathbb{R}^{n \times n}$: A real square matrix representing a linear transformation from $\mathbb{R}^n$ to $\mathbb{R}^n$.
    \item $\lambda \in \mathbb{C}$: An eigenvalue of $A$, satisfying $A\mathbf{v} = \lambda \mathbf{v}$.
    \item $\mathbf{v} \in \mathbb{C}^n \setminus \{\mathbf{0}\}$: A non-zero eigenvector associated with eigenvalue $\lambda$.
    \item $E_\lambda \triangleq \Null(A - \lambda I)$: The eigenspace corresponding to eigenvalue $\lambda$.
    \item $p_A(\lambda) \triangleq \det(A - \lambda I)$: The degree-$n$ characteristic polynomial of $A$.
    \item $am(\lambda)$: Algebraic multiplicity (the multiplicity of $\lambda$ as a root of $p_A(\lambda)$).
    \item $gm(\lambda) \triangleq \dim(E_\lambda) = n - \rank(A - \lambda I)$: Geometric multiplicity (the number of linearly independent eigenvectors for $\lambda$).
    \item $P = [\mathbf{v}_1, \dots, \mathbf{v}_n] \in \mathbb{C}^{n \times n}$: The modal matrix whose columns are linearly independent eigenvectors.
    \item $D = \diag(\lambda_1, \dots, \lambda_n)$: The diagonal spectral matrix of eigenvalues.
    \item $\Phi \in \mathbb{R}^{N \times (p+1)}$: The Vandermonde design matrix for degree-$p$ polynomial regression.
\end{itemize}

% ----------------------------------------------------------------------
% PART 2: THE GRAND PARADOX & MOTIVATION
% ----------------------------------------------------------------------
\section{The Grand Paradox: The Linear Model That Curves}

\subsection{The Linearity Illusion}
When newcomers encounter the phrase \textbf{Linear Regression}, they almost universally imagine fitting a straight line $y = mx + c$ through data points on a 2D plane. When presented with a parabolic dataset—such as projectile motion under gravity $y(t) = -\frac{1}{2}gt^2 + v_0 t + y_0$—the novice assumes that linear regression has failed, and that one must abandon linear models in favor of "non-linear algorithms."

\begin{quote}
\textit{\textbf{The Paradox:}} How can a strictly \textbf{linear} model fit a parabolas, cubics, or oscillatory curves? If an algorithm can fit an arbitrary polynomial curve, why is it still classified as "linear"?
\end{quote}

The resolution hinges on the distinction between \textbf{linearity in variables} versus \textbf{linearity in parameters}:
\begin{equation}
f(x; \mathbf{w}) = w_0 + w_1 x + w_2 x^2 + \dots + w_p x^p = \sum_{j=0}^p w_j \phi_j(x) = \mathbf{w}^T \boldsymbol{\phi}(x)
\end{equation}
While $f(x)$ is non-linear with respect to the input variable $x$, it is \textbf{strictly linear with respect to the parameter vector $\mathbf{w}$}:
\begin{equation}
f(x; \alpha \mathbf{u} + \beta \mathbf{v}) = \alpha f(x; \mathbf{u}) + \beta f(x; \mathbf{v})
\end{equation}
By defining a feature expansion $\boldsymbol{\phi}: \mathbb{R} \to \mathbb{R}^{p+1}$ where $\boldsymbol{\phi}(x) = [1, x, x^2, \dots, x^p]^T$, we transform the 1D non-linear curve into a \textbf{flat $(p+1)$-dimensional hyperplane} in feature space! Because the loss function $\|\mathbf{y} - \Phi \mathbf{w}\|_2^2$ remains a convex quadratic with respect to $\mathbf{w}$, all the guarantees of Ordinary Least Squares derived in Chapter~\ref{chap:week03} apply without alteration.

\begin{ponderbox}[The Non-Linear Illusion]
Consider the model $y = w_0 + w_1 \sin(x) + w_2 \exp(x) + w_3 \log(x)$. 
\begin{enumerate}
    \item Is this model linear regression or non-linear regression?
    \item What is the feature mapping $\boldsymbol{\phi}(x)$ that allows this model to be solved using the standard Normal Equations $\Phi^T \Phi \mathbf{w} = \Phi^T \mathbf{y}$?
    \item In contrast, consider the model $y = \frac{w_0}{1 + \exp(-w_1 x)}$. Can this model be solved via standard linear least squares? Why or why not?
\end{enumerate}
\end{ponderbox}

\subsection{The Coordinate Untangling Paradox}
Consider a general $n \times n$ matrix $A$. If you multiply an arbitrary vector $\mathbf{x} \in \mathbb{R}^n$ by $A$, the resulting vector $A\mathbf{x}$ typically rotates, skews, and stretches simultaneously. If we iterate this operation $k$ times:
\begin{equation}
\mathbf{x}_k = A \mathbf{x}_{k-1} = A^2 \mathbf{x}_{k-2} = \dots = A^k \mathbf{x}_0
\end{equation}
Every coordinate of $\mathbf{x}_k$ becomes an entangled combination of all previous coordinates.

Yet, in almost every linear system, there exist special directions where matrix multiplication does \textbf{not} rotate the vector at all! Along these privileged directions, matrix multiplication reduces to \textbf{pure scalar multiplication}:
\begin{equation}
A\mathbf{v} = \lambda \mathbf{v}
\end{equation}
These vectors $\mathbf{v}$ are the \textbf{eigenvectors}, and the scalars $\lambda$ are the \textbf{eigenvalues}. If we change coordinates to a basis of eigenvectors, the matrix $A$ becomes a simple diagonal matrix $D = \diag(\lambda_1, \dots, \lambda_n)$, and computing $A^k$ collapses to raising individual scalar numbers to the power $k$: $\lambda_i^k$.

% ----------------------------------------------------------------------
% PART 3: POLYNOMIAL REGRESSION & VANDERMONDE MATRICES
% ----------------------------------------------------------------------
\section{Polynomial Regression \& The Vandermonde Matrix}

Let $\{ (x_i, y_i) \}_{i=1}^N$ be a dataset of $N$ scalar observations. We wish to fit a degree-$p$ polynomial:
\begin{equation}
\hat{y}(x) = w_0 + w_1 x + w_2 x^2 + \dots + w_p x^p
\end{equation}
We define the polynomial feature mapping $\boldsymbol{\phi}: \mathbb{R} \to \mathbb{R}^{p+1}$:
\begin{equation}
\boldsymbol{\phi}(x) \triangleq \begin{bmatrix} 1 & x & x^2 & \dots & x^p \end{bmatrix}^T
\end{equation}
Stacking the feature vectors for all $N$ data points into rows forms the \textbf{Vandermonde Design Matrix}:
\begin{equation}
\Phi = \begin{bmatrix}
\boldsymbol{\phi}(x_1)^T \\
\boldsymbol{\phi}(x_2)^T \\
\vdots \\
\boldsymbol{\phi}(x_N)^T
\end{bmatrix} = \begin{bmatrix}
1 & x_1 & x_1^2 & \dots & x_1^p \\
1 & x_2 & x_2^2 & \dots & x_2^p \\
\vdots & \vdots & \vdots & \ddots & \vdots \\
1 & x_N & x_N^2 & \dots & x_N^p
\end{bmatrix} \in \mathbb{R}^{N \times (p+1)}
\label{eq:vandermonde_matrix}
\end{equation}
The polynomial regression problem is now formulated as the standard linear system:
\begin{equation}
\Phi \mathbf{w} \approx \mathbf{y} \quad \text{where } \mathbf{w} = \begin{bmatrix} w_0 \\ w_1 \\ \vdots \\ w_p \end{bmatrix} \in \mathbb{R}^{p+1}, \quad \mathbf{y} = \begin{bmatrix} y_1 \\ y_2 \\ \vdots \\ y_N \end{bmatrix} \in \mathbb{R}^N
\end{equation}

\begin{theorem}[Full Rank of Vandermonde Matrices]
Let $x_1, x_2, \dots, x_N$ be $N$ distinct real numbers. If $N \ge p + 1$, the Vandermonde matrix $\Phi \in \mathbb{R}^{N \times (p+1)}$ has full column rank: $\rank(\Phi) = p + 1$.
\end{theorem}

\begin{proof}
Suppose $\Phi \mathbf{c} = \mathbf{0}$ for some coefficient vector $\mathbf{c} = [c_0, c_1, \dots, c_p]^T \in \mathbb{R}^{p+1}$.
Each row $i$ of the product $\Phi \mathbf{c} = \mathbf{0}$ represents the evaluation of the polynomial $P(x) = \sum_{j=0}^p c_j x^j$ at point $x_i$:
\begin{equation}
P(x_i) = c_0 + c_1 x_i + c_2 x_i^2 + \dots + c_p x_i^p = 0 \quad \text{for all } i \in \{1, 2, \dots, N\}
\end{equation}
Thus, $P(x)$ is a polynomial of degree at most $p$ that has $N$ distinct roots $x_1, \dots, x_N$.
By the Fundamental Theorem of Algebra, a non-zero polynomial of degree $p$ can have at most $p$ distinct roots. Since $N \ge p + 1 > p$, the polynomial must be identically zero: $c_0 = c_1 = \dots = c_p = 0$.
Hence, the only vector in $\Null(\Phi)$ is $\mathbf{c} = \mathbf{0}$, proving $\rank(\Phi) = p + 1$.
\end{proof}

\begin{examplebox}{Polynomial Feature Expansion \& Exact Parabolic Interpolation}
Given three data points:
\begin{equation}
(x_1, y_1) = (-1, 1), \quad (x_2, y_2) = (0, 0), \quad (x_3, y_3) = (1, 2)
\end{equation}
Fit a quadratic polynomial $y = w_0 + w_1 x + w_2 x^2$ ($p = 2, N = 3$).

\textbf{Step 1: Construct the Square Vandermonde Matrix $\Phi$.}
\begin{equation}
\Phi = \begin{bmatrix} 1 & x_1 & x_1^2 \\ 1 & x_2 & x_2^2 \\ 1 & x_3 & x_3^2 \end{bmatrix} = \begin{bmatrix} 1 & -1 & (-1)^2 \\ 1 & 0 & 0^2 \\ 1 & 1 & 1^2 \end{bmatrix} = \begin{bmatrix} 1 & -1 & 1 \\ 1 & 0 & 0 \\ 1 & 1 & 1 \end{bmatrix}
\end{equation}
The target vector is $\mathbf{y} = [1, 0, 2]^T$.

\textbf{Step 2: Solve the Linear System $\Phi \mathbf{w} = \mathbf{y}$.}
Because $N = p + 1 = 3$, the system is exactly determined. We can solve via row operations:
From row 2: $1 \cdot w_0 + 0 \cdot w_1 + 0 \cdot w_2 = 0 \implies w_0 = 0$.
Substitute $w_0 = 0$ into rows 1 and 3:
\begin{align}
-w_1 + w_2 &= 1 \\
w_1 + w_2 &= 2
\end{align}
Adding the two equations: $2w_2 = 3 \implies w_2 = \frac{3}{2} = 1.5$.
Subtracting the equations: $2w_1 = 1 \implies w_1 = \frac{1}{2} = 0.5$.

\textbf{Step 3: Verification and Out-of-Sample Evaluation.}
The fitted polynomial is:
\begin{equation}
\hat{y}(x) = 0 + \frac{1}{2}x + \frac{3}{2}x^2
\end{equation}
Checking training nodes:
\begin{align}
\hat{y}(-1) &= 0.5(-1) + 1.5(-1)^2 = -0.5 + 1.5 = 1.0 \quad \checkmark \\
\hat{y}(0) &= 0 + 0 = 0.0 \quad \checkmark \\
\hat{y}(1) &= 0.5(1) + 1.5(1)^2 = 0.5 + 1.5 = 2.0 \quad \checkmark
\end{align}
Evaluating out-of-sample at $x = 2$:
\begin{equation}
\hat{y}(2) = 0.5(2) + 1.5(4) = 1 + 6 = 7.0
\end{equation}
\end{examplebox}

% ----------------------------------------------------------------------
% PART 4: SPECTRAL THEORY: EIGENVALUES & EIGENSPACES
% ----------------------------------------------------------------------
\section{Spectral Theory: Eigenvalues, Eigenspaces \& Multiplicity}

\subsection{The Characteristic Polynomial}
To find the invariant directions of an $n \times n$ matrix $A$, we seek non-zero vectors $\mathbf{v}$ and scalars $\lambda$ satisfying:
\begin{equation}
A\mathbf{v} = \lambda \mathbf{v} \iff (A - \lambda I)\mathbf{v} = \mathbf{0} \quad (\mathbf{v} \ne \mathbf{0})
\label{eq:eigen_fundamental}
\end{equation}
Equation~\eqref{eq:eigen_fundamental} states that the matrix $(A - \lambda I)$ maps a non-zero vector to zero. Therefore, $(A - \lambda I)$ must have a non-trivial nullspace, which means it cannot be invertible:
\begin{equation}
\det(A - \lambda I) = 0
\label{eq:char_eq}
\end{equation}
Expanding this determinant yields the \textbf{characteristic polynomial} $p_A(\lambda)$ of degree $n$:
\begin{equation}
p_A(\lambda) = (-1)^n \lambda^n + (-1)^{n-1} \trace(A) \lambda^{n-1} + \dots + \det(A) = 0
\end{equation}

\subsection{Algebraic vs. Geometric Multiplicity}

\begin{definitionbox}{Algebraic and Geometric Multiplicity}
Let $\lambda_0$ be an eigenvalue of $A \in \mathbb{R}^{n \times n}$:
\begin{enumerate}
    \item \textbf{Algebraic Multiplicity} $am(\lambda_0)$: The power of $(\lambda - \lambda_0)$ in the factorization of the characteristic polynomial:
    \begin{equation}
    p_A(\lambda) = (-1)^n (\lambda - \lambda_0)^{am(\lambda_0)} \prod_{j} (\lambda - \lambda_j)^{am(\lambda_j)}
    \end{equation}
    \item \textbf{Geometric Multiplicity} $gm(\lambda_0)$: The dimension of the eigenspace $E_{\lambda_0} = \Null(A - \lambda_0 I)$:
    \begin{equation}
    gm(\lambda_0) \triangleq \dim(E_{\lambda_0}) = n - \rank(A - \lambda_0 I)
    \end{equation}
\end{enumerate}
\end{definitionbox}

\begin{theorem}[Fundamental Multiplicity Inequality]
For every eigenvalue $\lambda$ of an $n \times n$ matrix $A$:
\begin{equation}
1 \le gm(\lambda) \le am(\lambda) \le n
\label{eq:mult_inequality}
\end{equation}
\end{theorem}

\begin{proof}
Because $\lambda$ is an eigenvalue, there exists at least one non-zero eigenvector, so $\dim(\Null(A - \lambda I)) \ge 1$, proving $gm(\lambda) \ge 1$.

To prove $gm(\lambda) \le am(\lambda)$, let $k = gm(\lambda)$, and choose an orthonormal basis $\{\mathbf{v}_1, \dots, \mathbf{v}_k\}$ for $E_\lambda$. Extend this to an orthonormal basis $\{\mathbf{v}_1, \dots, \mathbf{v}_k, \mathbf{u}_1, \dots, \mathbf{u}_{n-k}\}$ for all of $\mathbb{C}^n$. Assemble these vectors as columns of a unitary matrix $Q = [V_1 \mid U]$.

Because $A\mathbf{v}_i = \lambda \mathbf{v}_i$ for each $i \in \{1, \dots, k\}$, we have $A V_1 = \lambda V_1$.
Transforming $A$ by $Q^H (\cdot) Q$:
\begin{equation}
Q^H A Q = \begin{bmatrix} V_1^H \\ U^H \end{bmatrix} A \begin{bmatrix} V_1 & U \end{bmatrix} = \begin{bmatrix} V_1^H A V_1 & V_1^H A U \\ U^H A V_1 & U^H A U \end{bmatrix} = \begin{bmatrix} \lambda I_k & B \\ 0 & C \end{bmatrix}
\end{equation}
The characteristic polynomial of $Q^H A Q$ (which equals $p_A(t)$ because similar matrices share characteristic polynomials) is:
\begin{equation}
\det(Q^H A Q - t I_n) = \det\begin{bmatrix} (\lambda - t)I_k & B \\ 0 & C - t I_{n-k} \end{bmatrix} = (\lambda - t)^k \det(C - t I_{n-k})
\end{equation}
Thus, $(\lambda - t)^k$ is a factor of $p_A(t)$, which proves that the algebraic multiplicity of $\lambda$ must be at least $k = gm(\lambda)$.
\end{proof}

\begin{ponderbox}[The Geometric Multiplicity Bound]
Suppose a $4 \times 4$ matrix has characteristic polynomial $p(\lambda) = (\lambda - 3)^3 (\lambda + 1)$.
\begin{enumerate}
    \item What is the algebraic multiplicity of $\lambda = 3$?
    \item What are the possible values for the geometric multiplicity $gm(3)$?
    \item What is the geometric multiplicity of $\lambda = -1$? Can it ever equal 2? Why or why not?
\end{enumerate}
\end{ponderbox}

\begin{examplebox}{Diagnosing a Defective Shear Matrix}
Let matrix $B \in \mathbb{R}^{2 \times 2}$ be the horizontal shear operator:
\begin{equation}
B = \begin{bmatrix} 3 & 1 \\ 0 & 3 \end{bmatrix}
\end{equation}
Compute its characteristic polynomial, eigenvalues, algebraic and geometric multiplicities, and explain why $B$ is not diagonalizable.

\textbf{Step 1: Compute Characteristic Polynomial.}
\begin{equation}
p_B(\lambda) = \det(B - \lambda I) = \det\begin{bmatrix} 3 - \lambda & 1 \\ 0 & 3 - \lambda \end{bmatrix} = (3 - \lambda)^2 = 0
\end{equation}
The only eigenvalue is $\lambda = 3$, with algebraic multiplicity $am(3) = 2$.

\textbf{Step 2: Compute Eigenspace $E_3 = \Null(B - 3I)$.}
\begin{equation}
B - 3I = \begin{bmatrix} 3 - 3 & 1 \\ 0 & 3 - 3 \end{bmatrix} = \begin{bmatrix} 0 & 1 \\ 0 & 0 \end{bmatrix}
\end{equation}
Setting $(B - 3I)\mathbf{v} = \mathbf{0}$:
\begin{equation}
\begin{bmatrix} 0 & 1 \\ 0 & 0 \end{bmatrix} \begin{bmatrix} v_1 \\ v_2 \end{bmatrix} = \begin{bmatrix} 0 \\ 0 \end{bmatrix} \implies 0 \cdot v_1 + 1 \cdot v_2 = 0 \implies v_2 = 0
\end{equation}
The first coordinate $v_1$ is completely free. Thus, any eigenvector must be of the form $\mathbf{v} = [v_1, 0]^T = v_1 [1, 0]^T$.
The eigenspace is:
\begin{equation}
E_3 = \Span\left\{ \begin{bmatrix} 1 \\ 0 \end{bmatrix} \right\} \implies gm(3) = \dim(E_3) = 1
\end{equation}

\textbf{Step 3: Diagnosis of Defectiveness.}
Because $gm(3) = 1 < am(3) = 2$, matrix $B$ is \textbf{defective}. It possesses only one linearly independent eigenvector in $\mathbb{R}^2$. It is physically impossible to construct an invertible $2 \times 2$ modal matrix $P = [\mathbf{v}_1, \mathbf{v}_2]$, so $B$ \textbf{cannot be diagonalized}!
\end{examplebox}

% ----------------------------------------------------------------------
% PART 5: MATRIX DIAGONALIZATION & DYNAMICAL SYSTEMS
% ----------------------------------------------------------------------
\section{Matrix Diagonalization \& Iterated Powers}

\begin{theorem}[The Diagonalization Theorem]
An $n \times n$ matrix $A$ is diagonalizable if and only if the sum of geometric multiplicities equals $n$:
\begin{equation}
\sum_{i} gm(\lambda_i) = n \iff gm(\lambda_i) = am(\lambda_i) \quad \text{for all } i
\end{equation}
When diagonalizable, $A = P D P^{-1}$ where:
\begin{equation}
P = \begin{bmatrix} \mathbf{v}_1 & \mathbf{v}_2 & \dots & \mathbf{v}_n \end{bmatrix}, \quad D = \begin{bmatrix} \lambda_1 & & \\ & \ddots & \\ & & \lambda_n \end{bmatrix}
\label{eq:diagonalization_decomp}
\end{equation}
\end{theorem}

\begin{figure}[htbp]
\centering
\begin{tikzpicture}[
    scale=0.9,
    font=\sffamily\footnotesize
]

% Left: Original Unit Circle
\draw[thick, fill=primarylight!30] (-4.5, 0) circle (1.8cm);
\draw[->, >=stealth, color=gray] (-6.8, 0) -- (-2.2, 0) node[right] {$x_1$};
\draw[->, >=stealth, color=gray] (-4.5, -2.3) -- (-4.5, 2.3) node[above] {$x_2$};

% Eigenvectors on unit circle
\draw[->, >=stealth, line width=1.6pt, color=tealborder] (-4.5, 0) -- (-3.0, 0.9) node[above right] {$\mathbf{v}_1$};
\draw[->, >=stealth, line width=1.6pt, color=coralborder] (-4.5, 0) -- (-5.4, 1.5) node[above left] {$\mathbf{v}_2$};
\node[below] at (-4.5, -2.4) {\textbf{Input Space: } Unit Sphere $\|\mathbf{x}\| = 1$};

% Transformation arrow
\draw[->, >=stealth, line width=2pt, color=primarydark] (-1.8, 0) -- (0.8, 0) node[midway, above] {$\mathbf{x} \mapsto A\mathbf{x}$};

% Right: Transformed Ellipse
\begin{scope}[shift={(4.0, 0)}, rotate=30]
    \draw[thick, fill=tealbg!40] (0, 0) ellipse (3.2cm and 1.2cm);
    % Stretched eigenvectors
    \draw[->, >=stealth, line width=1.8pt, color=tealborder] (0, 0) -- (3.2, 0) node[above right] {$A\mathbf{v}_1 = \lambda_1 \mathbf{v}_1$};
    \draw[->, >=stealth, line width=1.8pt, color=coralborder] (0, 0) -- (0, 1.2) node[above left] {$A\mathbf{v}_2 = \lambda_2 \mathbf{v}_2$};
\end{scope}
\node[below] at (4.0, -2.4) {\textbf{Output Space: } Ellipsoid with axes $\lambda_i \mathbf{v}_i$};

\end{tikzpicture}
\caption{The Geometry of Eigendecomposition: The linear transformation $A$ deforms the unit circle into an ellipse whose principal axes align with the eigenvectors $\mathbf{v}_i$ scaled by eigenvalues $\lambda_i$.}
\label{fig:eigen_geometry}
\end{figure}

\begin{examplebox}{Step-by-Step Matrix Eigendecomposition and Power $A^k$}
Let $A = \begin{bmatrix} 4 & 2 \\ 1 & 3 \end{bmatrix}$. 
\begin{enumerate}
    \item Find eigenvalues $\lambda_1, \lambda_2$.
    \item Find eigenvectors $\mathbf{v}_1, \mathbf{v}_2$ and construct $P$ and $P^{-1}$.
    \item Compute the closed-form formula for $A^k$ and evaluate $A^5$.
\end{enumerate}

\textbf{Step 1: Compute Eigenvalues.}
\begin{equation}
\det(A - \lambda I) = \det\begin{bmatrix} 4 - \lambda & 2 \\ 1 & 3 - \lambda \end{bmatrix} = (4 - \lambda)(3 - \lambda) - 2 = \lambda^2 - 7\lambda + 10 = (\lambda - 5)(\lambda - 2) = 0
\end{equation}
The eigenvalues are $\lambda_1 = 5$ and $\lambda_2 = 2$. Both are distinct, guaranteeing that $A$ is diagonalizable.

\textbf{Step 2: Compute Eigenvectors.}
For $\lambda_1 = 5$:
\begin{equation}
(A - 5I)\mathbf{v}_1 = \begin{bmatrix} -1 & 2 \\ 1 & -2 \end{bmatrix} \begin{bmatrix} v_{11} \\ v_{12} \end{bmatrix} = \begin{bmatrix} 0 \\ 0 \end{bmatrix} \implies -v_{11} + 2v_{12} = 0 \implies \mathbf{v}_1 = \begin{bmatrix} 2 \\ 1 \end{bmatrix}
\end{equation}
For $\lambda_2 = 2$:
\begin{equation}
(A - 2I)\mathbf{v}_2 = \begin{bmatrix} 2 & 2 \\ 1 & 1 \end{bmatrix} \begin{bmatrix} v_{21} \\ v_{22} \end{bmatrix} = \begin{bmatrix} 0 \\ 0 \end{bmatrix} \implies v_{21} + v_{22} = 0 \implies \mathbf{v}_2 = \begin{bmatrix} 1 \\ -1 \end{bmatrix}
\end{equation}

\textbf{Step 3: Construct Modal Matrix $P$ and Invert.}
\begin{equation}
P = \begin{bmatrix} \mathbf{v}_1 & \mathbf{v}_2 \end{bmatrix} = \begin{bmatrix} 2 & 1 \\ 1 & -1 \end{bmatrix}, \quad \det(P) = 2(-1) - 1(1) = -3
\end{equation}
\begin{equation}
P^{-1} = \frac{1}{-3} \begin{bmatrix} -1 & -1 \\ -1 & 2 \end{bmatrix} = \frac{1}{3} \begin{bmatrix} 1 & 1 \\ 1 & -2 \end{bmatrix}
\end{equation}

\textbf{Step 4: Compute General Matrix Power $A^k = P D^k P^{-1}$.}
\begin{align}
A^k &= \begin{bmatrix} 2 & 1 \\ 1 & -1 \end{bmatrix} \begin{bmatrix} 5^k & 0 \\ 0 & 2^k \end{bmatrix} \left( \frac{1}{3} \begin{bmatrix} 1 & 1 \\ 1 & -2 \end{bmatrix} \right) \\
&= \frac{1}{3} \begin{bmatrix} 2 \cdot 5^k & 2^k \\ 5^k & -2^k \end{bmatrix} \begin{bmatrix} 1 & 1 \\ 1 & -2 \end{bmatrix} \\
&= \frac{1}{3} \begin{bmatrix} 2 \cdot 5^k + 2^k & 2 \cdot 5^k - 2 \cdot 2^k \\ 5^k - 2^k & 5^k + 2 \cdot 2^k \end{bmatrix}
\end{align}

\textbf{Step 5: Specific Evaluation for $k = 5$.}
Using $5^5 = 3125$ and $2^5 = 32$:
\begin{align}
(A^5)_{11} &= \frac{2(3125) + 32}{3} = \frac{6250 + 32}{3} = \frac{6282}{3} = 2094 \\
(A^5)_{12} &= \frac{2(3125) - 2(32)}{3} = \frac{6250 - 64}{3} = \frac{6186}{3} = 2062 \\
(A^5)_{21} &= \frac{3125 - 32}{3} = \frac{3093}{3} = 1031 \\
(A^5)_{22} &= \frac{3125 + 2(32)}{3} = \frac{3125 + 64}{3} = \frac{3189}{3} = 1063
\end{align}
Thus:
\begin{equation}
A^5 = \begin{bmatrix} 2094 & 2062 \\ 1031 & 1063 \end{bmatrix}
\end{equation}
Notice how the dominant eigenvalue $\lambda_1 = 5$ completely dictates the asymptotic scale of $A^k$ as $k \to \infty$!
\end{examplebox}

\begin{ponderbox}[The Spectral Radius and Dynamical Stability]
Let $\mathbf{x}_{t+1} = W \mathbf{x}_t$ describe the hidden state evolution of a linear recurrent neural network (RNN). The \textbf{spectral radius} is $\rho(W) \triangleq \max_i |\lambda_i|$.
\begin{enumerate}
    \item If $\rho(W) > 1$, what happens to $\|\mathbf{x}_t\|_2$ as $t \to \infty$? Why does this cause \textbf{exploding activations}?
    \item If $\rho(W) < 1$, what happens to $\|\mathbf{x}_t\|_2$ as $t \to \infty$? Why does this cause \textbf{vanishing gradients}?
    \item Why do modern recurrent architectures (e.g. orthogonal RNNs, unitary RNNs) initialize weights so that $\rho(W) \approx 1$?
\end{enumerate}
\end{ponderbox}

\subsection{Solving Discrete Dynamical Systems: Derivation of Binet's Fibonacci Formula}
The Fibonacci recurrence is defined as $F_{k+2} = F_{k+1} + F_k$ with initial conditions $F_0 = 0, F_1 = 1$.
We can express this second-order recurrence as a first-order two-dimensional linear dynamical system:
\begin{equation}
\mathbf{u}_{k+1} \triangleq \begin{bmatrix} F_{k+1} \\ F_k \end{bmatrix} = \begin{bmatrix} 1 & 1 \\ 1 & 0 \end{bmatrix} \begin{bmatrix} F_k \\ F_{k-1} \end{bmatrix} = A \mathbf{u}_k
\label{eq:fib_dyn}
\end{equation}
Iterating $k$ times from initial state $\mathbf{u}_0 = \begin{bmatrix} 1 \\ 0 \end{bmatrix}$:
\begin{equation}
\mathbf{u}_k = A^k \mathbf{u}_0 = A^k \begin{bmatrix} 1 \\ 0 \end{bmatrix}
\end{equation}
To compute $A^k$, we diagonalize $A = \begin{bmatrix} 1 & 1 \\ 1 & 0 \end{bmatrix}$:
\begin{equation}
\det(A - \lambda I) = \det\begin{bmatrix} 1 - \lambda & 1 \\ 1 & -\lambda \end{bmatrix} = \lambda^2 - \lambda - 1 = 0
\end{equation}
The eigenvalues are the Golden Ratio $\phi$ and its conjugate:
\begin{equation}
\lambda_1 = \frac{1 + \sqrt{5}}{2} = \phi \approx 1.618, \quad \lambda_2 = \frac{1 - \sqrt{5}}{2} = 1 - \phi = -\frac{1}{\phi} \approx -0.618
\end{equation}
Finding eigenvectors: $(A - \lambda_i I)\mathbf{v}_i = \mathbf{0} \implies \mathbf{v}_i = \begin{bmatrix} \lambda_i \\ 1 \end{bmatrix}$.
Thus:
\begin{equation}
P = \begin{bmatrix} \lambda_1 & \lambda_2 \\ 1 & 1 \end{bmatrix}, \quad P^{-1} = \frac{1}{\lambda_1 - \lambda_2} \begin{bmatrix} 1 & -\lambda_2 \\ -1 & \lambda_1 \end{bmatrix} = \frac{1}{\sqrt{5}} \begin{bmatrix} 1 & -\lambda_2 \\ -1 & \lambda_1 \end{bmatrix}
\end{equation}
Now, expanding $\mathbf{u}_k = P D^k P^{-1} \mathbf{u}_0$:
\begin{align}
\begin{bmatrix} F_{k+1} \\ F_k \end{bmatrix} &= \frac{1}{\sqrt{5}} \begin{bmatrix} \lambda_1 & \lambda_2 \\ 1 & 1 \end{bmatrix} \begin{bmatrix} \lambda_1^k & 0 \\ 0 & \lambda_2^k \end{bmatrix} \begin{bmatrix} 1 & -\lambda_2 \\ -1 & \lambda_1 \end{bmatrix} \begin{bmatrix} 1 \\ 0 \end{bmatrix} \nonumber \\
&= \frac{1}{\sqrt{5}} \begin{bmatrix} \lambda_1 & \lambda_2 \\ 1 & 1 \end{bmatrix} \begin{bmatrix} \lambda_1^k \\ -\lambda_2^k \end{bmatrix} = \frac{1}{\sqrt{5}} \begin{bmatrix} \lambda_1^{k+1} - \lambda_2^{k+1} \\ \lambda_1^k - \lambda_2^k \end{bmatrix}
\end{align}
Reading the second component yields Binet's celebrated formula:
\begin{equation}
F_k = \frac{1}{\sqrt{5}} \left( \left(\frac{1 + \sqrt{5}}{2}\right)^k - \left(\frac{1 - \sqrt{5}}{2}\right)^k \right)
\label{eq:binet}
\end{equation}

% ----------------------------------------------------------------------
% PART 6: EXPLICIT GOTCHAS & DIAGNOSTIC MISCONCEPTIONS
% ----------------------------------------------------------------------
\section{Explicit Gotchas \& Diagnostic Misconceptions}

\begin{trapbox}{Gotcha 1: Assuming Every Square Matrix is Diagonalizable}
A practitioner encounters matrix $J = \begin{bmatrix} 3 & 1 \\ 0 & 3 \end{bmatrix}$ and attempts to compute $J^{100} = P D^{100} P^{-1}$.
\textbf{The Trap:} The characteristic polynomial is $(\lambda - 3)^2 = 0$, giving eigenvalue $\lambda = 3$ with algebraic multiplicity $am(3) = 2$. However, the eigenspace is $\Null(J - 3I) = \Null\begin{bmatrix} 0 & 1 \\ 0 & 0 \end{bmatrix} = \Span\left\{\begin{bmatrix} 1 \\ 0 \end{bmatrix}\right\}$, which has dimension $gm(3) = 1$. Because $gm(3) < am(3)$, the matrix is \textbf{defective} (non-diagonalizable) and lacks a complete basis of eigenvectors. To compute powers of defective matrices, one must use the Jordan Normal Form.
\end{trapbox}

\begin{trapbox}{Gotcha 2: High-Degree Polynomial Regression and Runge's Phenomenon}
To achieve zero training error on $N = 15$ data points, a student fits a degree-$14$ polynomial, observing wild oscillations near the boundary edges of the interval $[-1, 1]$.
\textbf{The Trap:} This is \textbf{Runge's phenomenon}. In equispaced polynomial interpolation, the condition number of the Vandermonde matrix grows exponentially: $\kappa(\Phi) \sim \mathcal{O}(e^p)$. The resulting high-degree polynomial exhibits catastrophic out-of-sample test error. In production machine learning, high-degree polynomial expansions must be paired with Ridge regularization ($\lambda \|\mathbf{w}\|_2^2$), Chebyshev-spaced nodes, or replaced by local splines.
\end{trapbox}

\begin{trapbox}{Gotcha 3: Believing Real Matrices Always Have Real Eigenvalues}
A programmer computes the eigenvalues of a $2 \times 2$ 90-degree counterclockwise rotation matrix $R = \begin{bmatrix} 0 & -1 \\ 1 & 0 \end{bmatrix}$ and expects real scale factors.
\textbf{The Trap:} The characteristic equation is $\det(R - \lambda I) = \lambda^2 + 1 = 0 \implies \lambda = \pm i$. A pure rotation leaves no 1-dimensional real line invariant. Eigenvalues of real matrices are generally \textbf{complex numbers}. Real eigenvalues are guaranteed if and only if the matrix is symmetric ($A = A^T$), as established by the Spectral Theorem (Chapter~\ref{chap:week06}).
\end{trapbox}

\begin{trapbox}{Gotcha 4: Illegally Adding Eigenvalues for Matrix Sums}
An engineer knows the eigenvalues of $A$ are $\{2, 5\}$ and the eigenvalues of $B$ are $\{1, 3\}$. They conclude the eigenvalues of $A + B$ are $\{3, 8\}$.
\textbf{The Trap:} In general, $\lambda_i(A + B) \ne \lambda_i(A) + \lambda_i(B)$. The eigenvalues of a sum equal the sum of eigenvalues if and only if $A$ and $B$ share a simultaneous basis of eigenvectors (which occurs $\iff AB = BA$, meaning the matrices commute). For non-commuting matrices, the eigenvectors mix non-linearly.
\end{trapbox}

% ----------------------------------------------------------------------
% PART 7: COMPARATIVE SYNTHESIS & SUMMARY TABLE
% ----------------------------------------------------------------------
\section{Comparative Synthesis Table}

Table~\ref{tab:matrix_classes} compares the spectral properties across canonical matrix classes.

\begin{table}[htbp]
\centering
\small
\begin{tabularx}{\textwidth}{lXXXX}
\toprule
\textbf{Matrix Class} & \textbf{Symmetry Condition} & \textbf{Eigenvalue Field} & \textbf{Eigenvectors} & \textbf{Diagonalizability} \\
\midrule
General Square & $A \in \mathbb{R}^{n \times n}$ & $\lambda \in \mathbb{C}$ & Linearly independent if distinct & Diagonalizable $\iff gm(\lambda) = am(\lambda)$ \\
\addlinespace
Symmetric & $A = A^T$ & $\lambda \in \mathbb{R}$ strictly & Mutually orthogonal ($Q^T Q = I$) & Always orthogonally diagonalizable ($Q \Lambda Q^T$) \\
\addlinespace
Orthogonal & $Q^T Q = I$ & $|\lambda| = 1$ ($\lambda = e^{i\theta}$) & Mutually orthogonal in $\mathbb{C}^n$ & Always diagonalizable over $\mathbb{C}$ \\
\addlinespace
Projection & $P^2 = P = P^T$ & $\lambda \in \{0, 1\}$ strictly & Orthogonal bases for $V$ and $V^\perp$ & Always diagonalizable; $\rank(P) = \trace(P)$ \\
\addlinespace
Defective & $gm(\lambda) < am(\lambda)$ & $\lambda \in \mathbb{C}$ & Incomplete spanning set ($< n$) & Not diagonalizable; requires Jordan Normal Form \\
\bottomrule
\end{tabularx}
\caption{Comparative Synthesis of Spectral Properties Across Matrix Classes.}
\label{tab:matrix_classes}
\end{table}

% ----------------------------------------------------------------------
% PART 8: PYTHON / SCIENTIFIC COMPUTING LAB
% ----------------------------------------------------------------------
\section{Python / Scientific Computing Lab}

The production script below implements polynomial regression, evaluates the condition number explosion of Vandermonde matrices, and verifies Binet's Fibonacci formula via matrix diagonalization.

\begin{lstlisting}[language=Python, caption={Polynomial Regression and Spectral Diagonalization in Python.}]
import numpy as np

def fit_polynomial(x: np.ndarray, y: np.ndarray, degree: int):
    """Fits polynomial regression via Vandermonde expansion and evaluates condition number."""
    Phi = np.vander(x, N=degree + 1, increasing=True)
    cond_num = np.linalg.cond(Phi)
    # Solve via QR for numerical stability
    Q, R = np.linalg.qr(Phi)
    w = np.linalg.solve(R, Q.T @ y)
    return w, cond_num

def fibonacci_via_diagonalization(k: int) -> int:
    """Computes k-th Fibonacci number via spectral eigendecomposition of transition matrix."""
    A = np.array([[1.0, 1.0], [1.0, 0.0]])
    eigvals, P = np.linalg.eig(A)

    # Sort eigenvalues descending so lambda_1 is the Golden Ratio
    idx = np.argsort(eigvals)[::-1]
    eigvals = eigvals[idx]
    P = P[:, idx]

    # Compute A^k = P @ diag(lambda^k) @ P^{-1}
    D_k = np.diag(eigvals ** k)
    P_inv = np.linalg.inv(P)
    A_k = P @ D_k @ P_inv

    u0 = np.array([1.0, 0.0])
    u_k = A_k @ u0
    return int(np.round(u_k[1]))

def main():
    # 1. Vandermonde condition number explosion
    x_sample = np.linspace(-1.0, 1.0, 20)
    print("Vandermonde Matrix Condition Number vs. Degree:")
    for deg in [2, 5, 10, 15]:
        Phi = np.vander(x_sample, N=deg + 1, increasing=True)
        print(f"  Degree {deg:2d}: kappa(Phi) = {np.linalg.cond(Phi):.2e}")

    # 2. Fibonacci closed-form verification
    print("\nFibonacci Sequence via Matrix Powers:")
    for k in range(1, 11):
        f_k = fibonacci_via_diagonalization(k)
        print(f"  F_{k:02d} = {f_k}")

if __name__ == "__main__":
    main()
\end{lstlisting}

% ----------------------------------------------------------------------
% PART 9: COMPREHENSIVE PRACTICE PROBLEMS
% ----------------------------------------------------------------------
\section{Comprehensive Practice Problems}

\subsection*{Part A: Analytical \& Mathematical Problems}

\begin{enumerate}[label=\textbf{P\arabic*.}]
    \item \textbf{Eigendecomposition of a $2 \times 2$ Matrix:}
    Let $A = \begin{bmatrix} 4 & 2 \\ 1 & 3 \end{bmatrix}$.
    \begin{enumerate}[label=(\alph*)]
        \item Compute the characteristic equation $p_A(\lambda) = \det(A - \lambda I) = 0$ and find all eigenvalues.
        \item For each eigenvalue, find a corresponding eigenvector.
        \item Construct the modal matrix $P$ and diagonal matrix $D$ such that $A = PDP^{-1}$.
        \item Use this diagonalization to compute the exact matrix $A^5$.
    \end{enumerate}

    \item \textbf{Defective Matrix and Algebraic vs. Geometric Multiplicity:}
    Let $B = \begin{bmatrix} 2 & 1 & 0 \\ 0 & 2 & 0 \\ 0 & 0 & 5 \end{bmatrix}$.
    \begin{enumerate}[label=(\alph*)]
        \item Find the eigenvalues of $B$ and their algebraic multiplicities.
        \item Compute the eigenspace $E_\lambda = \Null(B - \lambda I)$ for each eigenvalue.
        \item Determine the geometric multiplicity of each eigenvalue.
        \item State whether $B$ is diagonalizable, justifying your answer rigorously.
    \end{enumerate}

    \item \textbf{Spectral Theorem for Symmetric Matrices:}
    Let $S \in \mathbb{R}^{n \times n}$ be a real symmetric matrix ($S = S^T$).
    \begin{enumerate}[label=(\alph*)]
        \item Prove that all eigenvalues of $S$ are real ($\lambda \in \mathbb{R}$).
        \item Prove that eigenvectors corresponding to distinct eigenvalues are mutually orthogonal: if $\lambda_1 \ne \lambda_2$, then $\mathbf{v}_1^T \mathbf{v}_2 = 0$.
    \end{enumerate}

    \item \textbf{Cayley-Hamilton Theorem Application:}
    The Cayley-Hamilton Theorem states that every square matrix satisfies its own characteristic equation: $p_A(A) = 0$.
    For matrix $A = \begin{bmatrix} 2 & 1 \\ 3 & 4 \end{bmatrix}$:
    \begin{enumerate}[label=(\alph*)]
        \item Find the characteristic polynomial $p_A(\lambda)$.
        \item Verify directly that $A^2 - \trace(A)A + \det(A)I = 0$.
        \item Use the Cayley-Hamilton equation to express $A^{-1}$ as a linear combination of $A$ and $I$.
    \end{enumerate}

    \item \textbf{Markov Transition Matrix Dynamics:}
    Consider a 2-state Markov transition matrix $T = \begin{bmatrix} 0.8 & 0.3 \\ 0.2 & 0.7 \end{bmatrix}$.
    \begin{enumerate}[label=(\alph*)]
        \item Prove that $\lambda_1 = 1$ is always an eigenvalue for any column-stochastic matrix.
        \item Find the second eigenvalue $\lambda_2$.
        \item Diagonalize $T = P D P^{-1}$ and compute the steady-state distribution $\boldsymbol{\pi} = \lim_{k \to \infty} T^k \mathbf{p}_0$ for an arbitrary initial probability vector $\mathbf{p}_0$.
    \end{enumerate}
\end{enumerate}

\subsection*{Part B: Diagnostic Multiple-Choice Questions}

\begin{enumerate}[label=\textbf{MCQ\arabic*.}]
    \item A matrix $A \in \mathbb{R}^{3 \times 3}$ has trace $\trace(A) = 6$ and determinant $\det(A) = 8$. Two of its eigenvalues are $\lambda_1 = 2$ and $\lambda_2 = 4$. What is the third eigenvalue $\lambda_3$, and is $A$ invertible?
    \begin{enumerate}[label=(\alph*)]
        \item $\lambda_3 = 0$; $A$ is singular.
        \item $\lambda_3 = 1$; $A$ is invertible.
        \item $\lambda_3 = 0$; $A$ is invertible.
        \item $\lambda_3 = -2$; $A$ is invertible.
    \end{enumerate}

    \item Why does fitting a degree-$p$ polynomial with $p = 15$ to $N = 16$ data points via Ordinary Least Squares typically perform poorly on new test data, even when training error is zero?
    \begin{enumerate}[label=(\alph*)]
        \item Because the Vandermonde matrix is non-invertible for $p < N$.
        \item Because high-degree polynomials have high bias and cannot fit the training points accurately.
        \item Because the model has excessive variance and overfits to training noise, exhibiting Runge's phenomenon at interval boundaries.
        \item Because polynomial regression requires the response variable $y$ to follow a Gaussian distribution.
    \end{enumerate}
\end{enumerate}
